
\documentclass{article}
\usepackage{colortbl}
\usepackage{makecell}
\usepackage{multirow}
\usepackage{supertabular}

\begin{document}

\newcounter{utterance}

\centering \large Interaction Transcript for game `dond', experiment `coop\_it', episode 2 with llama{-}3.3{-}70b{-}instruct{-}t1.0.
\vspace{24pt}

{ \footnotesize  \setcounter{utterance}{1}
\setlength{\tabcolsep}{0pt}
\begin{supertabular}{c@{$\;$}|p{.15\linewidth}@{}p{.15\linewidth}p{.15\linewidth}p{.15\linewidth}p{.15\linewidth}p{.15\linewidth}}
    \# & $\;$A & \multicolumn{4}{c}{Game Master} & $\;\:$B\\
    \hline

    \theutterance \stepcounter{utterance}  
    & & \multicolumn{4}{p{0.6\linewidth}}{
        \cellcolor[rgb]{0.9,0.9,0.9}{
            \makecell[{{p{\linewidth}}}]{
                \texttt{\tiny{[P1$\langle$GM]}}
                \texttt{State giocando a un gioco di negoziazione cooperativa in cui dovete accordarvi su come dividere una serie di oggetti tra voi e un altro giocatore.} \\
\\ 
\texttt{Regole:} \\
\texttt{(a) A Lei e all'altro giocatore viene dato un insieme di oggetti. Ognuno di voi riceve anche una funzione di valore segreta, che rappresenta il valore di ciascun tipo di oggetto.} \\
\texttt{(b) Si scambiano messaggi con l'altro giocatore per concordare chi si aggiudica gli oggetti. Potete inviare un massimo di 5 messaggi ciascuno, oppure terminare in anticipo facendo una proposta segreta in qualsiasi momento.} \\
\texttt{(c) A ciascuno di voi viene chiesto di inviare una proposta segreta indicando gli oggetti che desiderate, formattata tra parentesi quadre come segue: "{[}Proposta: <numero> <nome oggetto>, <numero> <nome oggetto>, <...>{]}".} \\
\texttt{(d) Se le vostre proposte sono complementari, cioè ci sono abbastanza oggetti per soddisfare entrambe le proposte, a ciascun giocatore viene assegnato un punteggio basato sulla somma dei valori degli oggetti ricevuti. In caso contrario, entrambi ricevono zero punti.} \\
\texttt{(e) L'obiettivo è massimizzare la somma del proprio punteggio e di quello dell'altro giocatore.} \\
\\ 
\texttt{Cominciamo.} \\
\\ 
\texttt{L'insieme degli oggetti disponibili è:} \\
\\ 
\texttt{1 telefono, 2 specchi, 2 frutti} \\
\\ 
\texttt{La funzione valore segreta è:} \\
\\ 
\texttt{telefono: 4, specchio: 2, frutta: 1} \\
\\ 
\texttt{IMPORTANTE: i vostri messaggi, a meno che non si tratti di una proposta segreta, vengono trasmessi direttamente all'altro giocatore, quindi non includete alcuna risposta alle regole o testo di annuncio del vostro messaggio. Per fare una proposta segreta, utilizzate il formato indicato. Non utilizzare delle parentesi quadre quando si comunica all'altro giocatore, altrimenti verrà interpretata come una proposta segreta.} \\
            }
        }
    }
    & \\ \\

    \theutterance \stepcounter{utterance}  
    & \multicolumn{4}{p{0.6\linewidth}}{
        \cellcolor[rgb]{0.9,0.9,0.9}{
            \makecell[{{p{\linewidth}}}]{
                \texttt{\tiny{[P1$\rangle$GM]}}
                \texttt{Ciao, vorrei iniziare la negoziazione. Il mio obiettivo è trovare un accordo che massimizzi la somma dei nostri punteggi. Quanti oggetti pensi di voler prendere?} \\
            }
        }
    }
    & & \\ \\

    \theutterance \stepcounter{utterance}  
    & & & \multicolumn{2}{p{0.3\linewidth}}{
        \cellcolor[rgb]{0.9,0.9,0.9}{
            \makecell[{{p{\linewidth}}}]{
                \texttt{\tiny{[GM$|$GM]}}
                \texttt{continue} \\
            }
        }
    }
    & & \\ \\

    \theutterance \stepcounter{utterance}  
    & & & \multicolumn{2}{p{0.3\linewidth}}{
        \cellcolor[rgb]{0.9,0.9,0.9}{
            \makecell[{{p{\linewidth}}}]{
                \texttt{\tiny{[GM$|$GM]}}
                \texttt{Ciao, vorrei iniziare la negoziazione. Il mio obiettivo è trovare un accordo che massimizzi la somma dei nostri punteggi. Quanti oggetti pensi di voler prendere?} \\
            }
        }
    }
    & & \\ \\

    \theutterance \stepcounter{utterance}  
    & & \multicolumn{4}{p{0.6\linewidth}}{
        \cellcolor[rgb]{0.9,0.9,0.9}{
            \makecell[{{p{\linewidth}}}]{
                \texttt{\tiny{[GM$\rangle$P2]}}
                \texttt{State giocando a un gioco di negoziazione cooperativa in cui dovete accordarvi su come dividere una serie di oggetti tra voi e un altro giocatore.} \\
\\ 
\texttt{Regole:} \\
\texttt{(a) A Lei e all'altro giocatore viene dato un insieme di oggetti. Ognuno di voi riceve anche una funzione di valore segreta, che rappresenta il valore di ciascun tipo di oggetto.} \\
\texttt{(b) Si scambiano messaggi con l'altro giocatore per concordare chi si aggiudica gli oggetti. Potete inviare un massimo di 5 messaggi ciascuno, oppure terminare in anticipo facendo una proposta segreta in qualsiasi momento.} \\
\texttt{(c) A ciascuno di voi viene chiesto di inviare una proposta segreta indicando gli oggetti che desiderate, formattata tra parentesi quadre come segue: "{[}Proposta: <numero> <nome oggetto>, <numero> <nome oggetto>, <...>{]}".} \\
\texttt{(d) Se le vostre proposte sono complementari, cioè ci sono abbastanza oggetti per soddisfare entrambe le proposte, a ciascun giocatore viene assegnato un punteggio basato sulla somma dei valori degli oggetti ricevuti. In caso contrario, entrambi ricevono zero punti.} \\
\texttt{(e) L'obiettivo è massimizzare la somma del proprio punteggio e di quello dell'altro giocatore.} \\
\\ 
\texttt{Cominciamo.} \\
\\ 
\texttt{L'insieme degli oggetti disponibili è:} \\
\\ 
\texttt{1 telefono, 2 specchi, 2 frutti} \\
\\ 
\texttt{La funzione valore segreta è:} \\
\\ 
\texttt{telefono: 6, specchio: 0, frutta: 2} \\
\\ 
\texttt{IMPORTANTE: i vostri messaggi, a meno che non si tratti di una proposta segreta, vengono trasmessi direttamente all'altro giocatore, quindi non includete alcuna risposta alle regole o testo di annuncio del vostro messaggio. Per fare una proposta segreta, utilizzate il formato indicato. Non utilizzare delle parentesi quadre quando si comunica all'altro giocatore, altrimenti verrà interpretata come una proposta segreta.} \\
\\ 
\\ 
\texttt{Ciao, vorrei iniziare la negoziazione. Il mio obiettivo è trovare un accordo che massimizzi la somma dei nostri punteggi. Quanti oggetti pensi di voler prendere?} \\
            }
        }
    }
    & \\ \\

    \theutterance \stepcounter{utterance}  
    & & & \multicolumn{4}{p{0.6\linewidth}}{
        \cellcolor[rgb]{0.9,0.9,0.9}{
            \makecell[{{p{\linewidth}}}]{
                \texttt{\tiny{[GM$\langle$P2]}}
                \texttt{Ciao! Sono contento di iniziare la negoziazione. Non ho una idea precisa sul numero di oggetti che voglio prendere, ma sono disposto a considerare varie opzioni. Mi sembra che ci siano alcuni oggetti che potrebbero essere più preziosi per uno di noi due rispetto all'altro. Vorrei sapere, hai una preferenza per uno degli oggetti in particolare?} \\
            }
        }
    }
     \\ \\

    \theutterance \stepcounter{utterance}  
    & & & \multicolumn{2}{p{0.3\linewidth}}{
        \cellcolor[rgb]{0.9,0.9,0.9}{
            \makecell[{{p{\linewidth}}}]{
                \texttt{\tiny{[GM$|$GM]}}
                \texttt{continue} \\
            }
        }
    }
    & & \\ \\

    \theutterance \stepcounter{utterance}  
    & & & \multicolumn{2}{p{0.3\linewidth}}{
        \cellcolor[rgb]{0.9,0.9,0.9}{
            \makecell[{{p{\linewidth}}}]{
                \texttt{\tiny{[GM$|$GM]}}
                \texttt{Ciao! Sono contento di iniziare la negoziazione. Non ho una idea precisa sul numero di oggetti che voglio prendere, ma sono disposto a considerare varie opzioni. Mi sembra che ci siano alcuni oggetti che potrebbero essere più preziosi per uno di noi due rispetto all'altro. Vorrei sapere, hai una preferenza per uno degli oggetti in particolare?} \\
            }
        }
    }
    & & \\ \\

    \theutterance \stepcounter{utterance}  
    & & \multicolumn{4}{p{0.6\linewidth}}{
        \cellcolor[rgb]{0.9,0.9,0.9}{
            \makecell[{{p{\linewidth}}}]{
                \texttt{\tiny{[P1$\langle$GM]}}
                \texttt{Ciao! Sono contento di iniziare la negoziazione. Non ho una idea precisa sul numero di oggetti che voglio prendere, ma sono disposto a considerare varie opzioni. Mi sembra che ci siano alcuni oggetti che potrebbero essere più preziosi per uno di noi due rispetto all'altro. Vorrei sapere, hai una preferenza per uno degli oggetti in particolare?} \\
            }
        }
    }
    & \\ \\

    \theutterance \stepcounter{utterance}  
    & \multicolumn{4}{p{0.6\linewidth}}{
        \cellcolor[rgb]{0.9,0.9,0.9}{
            \makecell[{{p{\linewidth}}}]{
                \texttt{\tiny{[P1$\rangle$GM]}}
                \texttt{Sì, credo che il telefono sia un oggetto molto prezioso, ma non so se vale la pena contenderselo. Gli specchi e la frutta potrebbero essere più facili da dividere. Che ne pensi degli specchi? Ti interesserebbero entrambi o solo uno?} \\
            }
        }
    }
    & & \\ \\

    \theutterance \stepcounter{utterance}  
    & & & \multicolumn{2}{p{0.3\linewidth}}{
        \cellcolor[rgb]{0.9,0.9,0.9}{
            \makecell[{{p{\linewidth}}}]{
                \texttt{\tiny{[GM$|$GM]}}
                \texttt{continue} \\
            }
        }
    }
    & & \\ \\

    \theutterance \stepcounter{utterance}  
    & & & \multicolumn{2}{p{0.3\linewidth}}{
        \cellcolor[rgb]{0.9,0.9,0.9}{
            \makecell[{{p{\linewidth}}}]{
                \texttt{\tiny{[GM$|$GM]}}
                \texttt{Sì, credo che il telefono sia un oggetto molto prezioso, ma non so se vale la pena contenderselo. Gli specchi e la frutta potrebbero essere più facili da dividere. Che ne pensi degli specchi? Ti interesserebbero entrambi o solo uno?} \\
            }
        }
    }
    & & \\ \\

    \theutterance \stepcounter{utterance}  
    & & \multicolumn{4}{p{0.6\linewidth}}{
        \cellcolor[rgb]{0.9,0.9,0.9}{
            \makecell[{{p{\linewidth}}}]{
                \texttt{\tiny{[GM$\rangle$P2]}}
                \texttt{Sì, credo che il telefono sia un oggetto molto prezioso, ma non so se vale la pena contenderselo. Gli specchi e la frutta potrebbero essere più facili da dividere. Che ne pensi degli specchi? Ti interesserebbero entrambi o solo uno?} \\
            }
        }
    }
    & \\ \\

    \theutterance \stepcounter{utterance}  
    & & & \multicolumn{4}{p{0.6\linewidth}}{
        \cellcolor[rgb]{0.9,0.9,0.9}{
            \makecell[{{p{\linewidth}}}]{
                \texttt{\tiny{[GM$\langle$P2]}}
                \texttt{Gli specchi potrebbero essere un buon punto di partenza per la negoziazione. Io non sono così interessato a entrambi gli specchi, uno potrebbe essere sufficiente. Ma non sono disposto a rinunciare alla possibilità di prendere anche l'altro se si può raggiungere un accordo vantaggioso per entrambi. La frutta potrebbe essere un buon oggetto da scambiare o condividere. Che ne pensi del telefono? Sei disposto a lasciarlo a me se io fossi disposto a cedere su altri oggetti?} \\
            }
        }
    }
     \\ \\

    \theutterance \stepcounter{utterance}  
    & & & \multicolumn{2}{p{0.3\linewidth}}{
        \cellcolor[rgb]{0.9,0.9,0.9}{
            \makecell[{{p{\linewidth}}}]{
                \texttt{\tiny{[GM$|$GM]}}
                \texttt{continue} \\
            }
        }
    }
    & & \\ \\

    \theutterance \stepcounter{utterance}  
    & & & \multicolumn{2}{p{0.3\linewidth}}{
        \cellcolor[rgb]{0.9,0.9,0.9}{
            \makecell[{{p{\linewidth}}}]{
                \texttt{\tiny{[GM$|$GM]}}
                \texttt{Gli specchi potrebbero essere un buon punto di partenza per la negoziazione. Io non sono così interessato a entrambi gli specchi, uno potrebbe essere sufficiente. Ma non sono disposto a rinunciare alla possibilità di prendere anche l'altro se si può raggiungere un accordo vantaggioso per entrambi. La frutta potrebbe essere un buon oggetto da scambiare o condividere. Che ne pensi del telefono? Sei disposto a lasciarlo a me se io fossi disposto a cedere su altri oggetti?} \\
            }
        }
    }
    & & \\ \\

    \theutterance \stepcounter{utterance}  
    & & \multicolumn{4}{p{0.6\linewidth}}{
        \cellcolor[rgb]{0.9,0.9,0.9}{
            \makecell[{{p{\linewidth}}}]{
                \texttt{\tiny{[P1$\langle$GM]}}
                \texttt{Gli specchi potrebbero essere un buon punto di partenza per la negoziazione. Io non sono così interessato a entrambi gli specchi, uno potrebbe essere sufficiente. Ma non sono disposto a rinunciare alla possibilità di prendere anche l'altro se si può raggiungere un accordo vantaggioso per entrambi. La frutta potrebbe essere un buon oggetto da scambiare o condividere. Che ne pensi del telefono? Sei disposto a lasciarlo a me se io fossi disposto a cedere su altri oggetti?} \\
            }
        }
    }
    & \\ \\

    \theutterance \stepcounter{utterance}  
    & \multicolumn{4}{p{0.6\linewidth}}{
        \cellcolor[rgb]{0.9,0.9,0.9}{
            \makecell[{{p{\linewidth}}}]{
                \texttt{\tiny{[P1$\rangle$GM]}}
                \texttt{Il telefono è un oggetto molto importante per me, ma sono disposto a considerare di lasciarlo a te se ottengo una buona offerta in cambio. Gli specchi e la frutta potrebbero essere una buona combinazione per me. Se tu prendessi il telefono e io prendessi entrambi gli specchi e forse una delle frutte, potrebbe essere un buon equilibrio. Ma dobbiamo anche considerare che la nostra valutazione degli oggetti potrebbe essere diversa, quindi è difficile dire cosa sia una "buona offerta" senza conoscere i nostri valori segreti. Cosa ne pensi di questa possibile divisione?} \\
            }
        }
    }
    & & \\ \\

    \theutterance \stepcounter{utterance}  
    & & & \multicolumn{2}{p{0.3\linewidth}}{
        \cellcolor[rgb]{0.9,0.9,0.9}{
            \makecell[{{p{\linewidth}}}]{
                \texttt{\tiny{[GM$|$GM]}}
                \texttt{continue} \\
            }
        }
    }
    & & \\ \\

    \theutterance \stepcounter{utterance}  
    & & & \multicolumn{2}{p{0.3\linewidth}}{
        \cellcolor[rgb]{0.9,0.9,0.9}{
            \makecell[{{p{\linewidth}}}]{
                \texttt{\tiny{[GM$|$GM]}}
                \texttt{Il telefono è un oggetto molto importante per me, ma sono disposto a considerare di lasciarlo a te se ottengo una buona offerta in cambio. Gli specchi e la frutta potrebbero essere una buona combinazione per me. Se tu prendessi il telefono e io prendessi entrambi gli specchi e forse una delle frutte, potrebbe essere un buon equilibrio. Ma dobbiamo anche considerare che la nostra valutazione degli oggetti potrebbe essere diversa, quindi è difficile dire cosa sia una "buona offerta" senza conoscere i nostri valori segreti. Cosa ne pensi di questa possibile divisione?} \\
            }
        }
    }
    & & \\ \\

    \theutterance \stepcounter{utterance}  
    & & \multicolumn{4}{p{0.6\linewidth}}{
        \cellcolor[rgb]{0.9,0.9,0.9}{
            \makecell[{{p{\linewidth}}}]{
                \texttt{\tiny{[GM$\rangle$P2]}}
                \texttt{Il telefono è un oggetto molto importante per me, ma sono disposto a considerare di lasciarlo a te se ottengo una buona offerta in cambio. Gli specchi e la frutta potrebbero essere una buona combinazione per me. Se tu prendessi il telefono e io prendessi entrambi gli specchi e forse una delle frutte, potrebbe essere un buon equilibrio. Ma dobbiamo anche considerare che la nostra valutazione degli oggetti potrebbe essere diversa, quindi è difficile dire cosa sia una "buona offerta" senza conoscere i nostri valori segreti. Cosa ne pensi di questa possibile divisione?} \\
            }
        }
    }
    & \\ \\

    \theutterance \stepcounter{utterance}  
    & & & \multicolumn{4}{p{0.6\linewidth}}{
        \cellcolor[rgb]{0.9,0.9,0.9}{
            \makecell[{{p{\linewidth}}}]{
                \texttt{\tiny{[GM$\langle$P2]}}
                \texttt{La proposta che hai fatto sembra ragionevole, ma come hai detto, la valutazione degli oggetti potrebbe essere diversa per ciascuno di noi. Io sono disposto a considerare la tua proposta, ma vorrei anche esplorare altre opzioni. Ad esempio, se io prendessi il telefono e uno degli specchi, e tu prendessi l'altro specchio e entrambe le frutte, potrebbe essere un'altra divisione equilibrata. Tuttavia, senza conoscere i valori segreti, è difficile dire se questa divisione sarebbe vantaggiosa per entrambi. Forse potremmo provare a fare una proposta segreta e vedere se le nostre richieste sono compatibili. Sei pronto a fare una proposta segreta?} \\
            }
        }
    }
     \\ \\

    \theutterance \stepcounter{utterance}  
    & & & \multicolumn{2}{p{0.3\linewidth}}{
        \cellcolor[rgb]{0.9,0.9,0.9}{
            \makecell[{{p{\linewidth}}}]{
                \texttt{\tiny{[GM$|$GM]}}
                \texttt{continue} \\
            }
        }
    }
    & & \\ \\

    \theutterance \stepcounter{utterance}  
    & & & \multicolumn{2}{p{0.3\linewidth}}{
        \cellcolor[rgb]{0.9,0.9,0.9}{
            \makecell[{{p{\linewidth}}}]{
                \texttt{\tiny{[GM$|$GM]}}
                \texttt{La proposta che hai fatto sembra ragionevole, ma come hai detto, la valutazione degli oggetti potrebbe essere diversa per ciascuno di noi. Io sono disposto a considerare la tua proposta, ma vorrei anche esplorare altre opzioni. Ad esempio, se io prendessi il telefono e uno degli specchi, e tu prendessi l'altro specchio e entrambe le frutte, potrebbe essere un'altra divisione equilibrata. Tuttavia, senza conoscere i valori segreti, è difficile dire se questa divisione sarebbe vantaggiosa per entrambi. Forse potremmo provare a fare una proposta segreta e vedere se le nostre richieste sono compatibili. Sei pronto a fare una proposta segreta?} \\
            }
        }
    }
    & & \\ \\

    \theutterance \stepcounter{utterance}  
    & & \multicolumn{4}{p{0.6\linewidth}}{
        \cellcolor[rgb]{0.9,0.9,0.9}{
            \makecell[{{p{\linewidth}}}]{
                \texttt{\tiny{[P1$\langle$GM]}}
                \texttt{La proposta che hai fatto sembra ragionevole, ma come hai detto, la valutazione degli oggetti potrebbe essere diversa per ciascuno di noi. Io sono disposto a considerare la tua proposta, ma vorrei anche esplorare altre opzioni. Ad esempio, se io prendessi il telefono e uno degli specchi, e tu prendessi l'altro specchio e entrambe le frutte, potrebbe essere un'altra divisione equilibrata. Tuttavia, senza conoscere i valori segreti, è difficile dire se questa divisione sarebbe vantaggiosa per entrambi. Forse potremmo provare a fare una proposta segreta e vedere se le nostre richieste sono compatibili. Sei pronto a fare una proposta segreta?} \\
            }
        }
    }
    & \\ \\

    \theutterance \stepcounter{utterance}  
    & \multicolumn{4}{p{0.6\linewidth}}{
        \cellcolor[rgb]{0.9,0.9,0.9}{
            \makecell[{{p{\linewidth}}}]{
                \texttt{\tiny{[P1$\rangle$GM]}}
                \texttt{Sì, credo che sia arrivato il momento di fare una proposta segreta. Spero che le nostre richieste siano compatibili e che possiamo trovare un accordo vantaggioso per entrambi. {[}Proposta: 1 telefono, 1 specchio, 1 frutta{]}} \\
            }
        }
    }
    & & \\ \\

    \theutterance \stepcounter{utterance}  
    & & & \multicolumn{2}{p{0.3\linewidth}}{
        \cellcolor[rgb]{0.9,0.9,0.9}{
            \makecell[{{p{\linewidth}}}]{
                \texttt{\tiny{[GM$|$GM]}}
                \texttt{proposal} \\
            }
        }
    }
    & & \\ \\

    \theutterance \stepcounter{utterance}  
    & & & \multicolumn{2}{p{0.3\linewidth}}{
        \cellcolor[rgb]{0.9,0.9,0.9}{
            \makecell[{{p{\linewidth}}}]{
                \texttt{\tiny{[GM$|$GM]}}
                [1, 1, 1]
            }
        }
    }
    & & \\ \\

    \theutterance \stepcounter{utterance}  
    & & \multicolumn{4}{p{0.6\linewidth}}{
        \cellcolor[rgb]{0.9,0.9,0.9}{
            \makecell[{{p{\linewidth}}}]{
                \texttt{\tiny{[GM$\rangle$P2]}}
                \texttt{L'altro giocatore ha presentato una proposta segreta. Ora tocca a Lei presentare una sua proposta.} \\
            }
        }
    }
    & \\ \\

    \theutterance \stepcounter{utterance}  
    & & & \multicolumn{4}{p{0.6\linewidth}}{
        \cellcolor[rgb]{0.9,0.9,0.9}{
            \makecell[{{p{\linewidth}}}]{
                \texttt{\tiny{[GM$\langle$P2]}}
                \texttt{{[}Proposta: 1 telefono, 0 specchi, 2 frutti{]}} \\
            }
        }
    }
     \\ \\

    \theutterance \stepcounter{utterance}  
    & & & \multicolumn{2}{p{0.3\linewidth}}{
        \cellcolor[rgb]{0.9,0.9,0.9}{
            \makecell[{{p{\linewidth}}}]{
                \texttt{\tiny{[GM$|$GM]}}
                \texttt{proposal} \\
            }
        }
    }
    & & \\ \\

    \theutterance \stepcounter{utterance}  
    & & & \multicolumn{2}{p{0.3\linewidth}}{
        \cellcolor[rgb]{0.9,0.9,0.9}{
            \makecell[{{p{\linewidth}}}]{
                \texttt{\tiny{[GM$|$GM]}}
                [1, 0, 2]
            }
        }
    }
    & & \\ \\

    \theutterance \stepcounter{utterance}  
    & & & \multicolumn{2}{p{0.3\linewidth}}{
        \cellcolor[rgb]{0.9,0.9,0.9}{
            \makecell[{{p{\linewidth}}}]{
                \texttt{\tiny{[GM$|$GM]}}
                [[1, 1, 1], [1, 0, 2]]
            }
        }
    }
    & & \\ \\

\end{supertabular}
}

\end{document}
